% ============================================================
%  chapters/taskbook_body.tex —— 本科毕业设计（论文）任务书正文
%  内容源：thesis/task-book.md（唯一内容来源）
%  参考文献：直录源稿文本（不用 .bib、不用 \cite），序号由 tjreflist 自动生成。
%  版式约定：latex/CONVENTIONS.md；正文分节全部交给 \tjformsection /
%            \tjformsubsection / \tjformsubsub，不手写 \heiti / \zihao。
%  结构：官方任务书表格的六个小节。源稿「一、课题基本信息」已由封面
%        （chapters/metadata.tex + \MakeTaskBook）自动排版，不在此重复。
%  篇幅：一至四部分合计约 1.6 页（实测 1.57 页）；课题背景不列参考文献。
%        全文 4 页（封面 + 3 页），在学院 3—7 页要求内。第 2 页已排满，
%        再增字会把「三、」的清单条目挤到第 3 页，故 v6 起正文按等字数改写。
%  口径：封面「共 32 周」由 metadata.tex 的 \taskbookperiod 按起讫含首尾
%        向上取整自动换算（2026-10-26—2027-05-31，218 天 = 31.14 周 → 32 周），
%        第 32 周为收尾周，不是笔误；不要手改周数，也不要改截止日（会变 33 周）。
% ============================================================


% ============================================================
\tjformsection{一、课题背景}\tjformsecnum{1}

代码智能体（code agent）已逐渐应用于日常开发：开发者以自然语言描述任务，由智能体完成代码分析、文件修改、工具调用、测试执行与问题排查。但对话、代码差异、工具调用与测试结果通常分散在不同界面。由此产生三方面问题：一是缺少任务视角，多轮修改、多任务并行与中断重启后难以掌握当前进展；二是修改原因难以回溯，代码因何而改、依据哪次测试验证，需依靠人工查找与回忆；三是中断后上下文丢失，重新开始需重读长对话与大量差异。

本课题面向单个开发者，研究开发任务、对话、代码变更、工具调用与测试结果的关联与组织方法，设计并实现支持任务上下文追踪与变更溯源的原型系统，帮助开发者围绕任务、文件或代码片段查看证据链，并在中断后恢复上下文。

% ============================================================
\tjformsection{二、应完成的主要内容}\tjformsecnum{2}

本课题要求建立开发任务、对话、工具调用、文件变更、测试运行等实体的信息模型，设计统一的事件表示、可扩展的事件模式与存储结构；同时接入至少一种真实代码智能体的事件源，结合 Git 与测试输出采集开发过程数据，研究将异构记录关联到开发任务的方法。

在此基础上，以变更块为粒度生成“变更卡片”，记录触发来源、实施过程、修改前后快照、验证依据与后续演化，支持“代码—原因—证据”的双向查询；实现任务中心、任务时间线与变更证据侧栏三个视图，任务中断后生成任务快照，支持由快照跳转到证据；同时探索利用大语言模型生成任务摘要、整理开发进展（输出须绑定原始证据，证据不足时弃答），并在典型编程任务上与基线方法对比，评估其在代码理解、修改原因分析、问题定位与任务衔接方面的作用。

% ============================================================
\tjformsection{三、基本要求及应完成的成果形式}\tjformsecnum{3}

本课题要求学生完成以下工作：

\begin{enumerate}[label=\arabic*.,leftmargin=2.2em,labelsep=0.4em]
  \item 调研代码智能体、开发过程追踪与变更溯源的研究与产品现状；
  \item 设计开发任务、对话、工具调用、文件变更、测试运行等实体的信息模型；
  \item 实现任务上下文追踪与变更溯源原型系统，接入真实事件源；
  \item 在典型编程任务上开展实验，与基线方法对比并形成可复现数据；
  \item 撰写毕业设计论文，说明需求分析、系统设计、实现与评估过程。
\end{enumerate}

\newpage
成果形式：可运行的原型系统一套（含源代码、部署与使用说明）；毕业设计（论文）一份；系统设计、测试与评估报告，以及实验数据；演示视频与答辩幻灯片。

% ============================================================
\tjformsection{四、进度安排}\tjformsecnum{4}

\begin{table}[H]
  \centering
  \caption{进度安排}\label{tab:tb-schedule}
  \begin{tabularx}{\linewidth}{@{}>{\centering\arraybackslash}p{3em}>{\centering\arraybackslash}X>{\centering\arraybackslash}p{10em}@{}}
    \toprule
    序号 & 设计（论文）各阶段名称 & 时间安排（教学周） \\
    \midrule
    1 & 文献研究与产品调研 & 1—8 \\
    2 & 关键方法设计与实现 & 9—14 \\
    3 & 系统设计与实现 & 15—20 \\
    4 & 系统验证与测试 & 21—22 \\
    5 & 毕业设计论文撰写 & 23—28 \\
    6 & 论文审阅与修改 & 29—32 \\
    \bottomrule
  \end{tabularx}
\end{table}

第 9 周进行开题；时间安排以教学周计，具体以学院通知为准。

% ============================================================
\tjformsection{五、主要参考文献}\tjformsecnum{5}

% 参考文献直录（不用 .bib、不用 \cite）：条目文本照抄源稿「七、主要参考文献」[1]—[18]，
% 含 [C]//、[J]、[EB/OL]、(2026-10-27) 等标记，不删不改；序号由 tjreflist 环境自动生成。
% 字号与行距照已审定的 Word 定稿（deliverables/毕业设计任务书-v1-20261027.docx：
% 文献条目为五号 10.5 pt、单倍行距），故此处按五号单倍行距排版，不随正文小四。
{\zihao{5}\linespread{1}\selectfont
\begin{tjreflist}
  \item JIMENEZ C E, YANG J, WETTIG A, et al. SWE-bench: can language models resolve real-world GitHub issues?[C]//Proceedings of the 12th International Conference on Learning Representations (ICLR 2024). 2024. arXiv: 2310.06770.
  \item YANG J, JIMENEZ C E, WETTIG A, et al. SWE-agent: agent-computer interfaces enable automated software engineering[C]//Advances in Neural Information Processing Systems 37 (NeurIPS 2024). 2024. arXiv: 2405.15793.
  \item WANG X, LI B, SONG Y, et al. OpenHands: an open platform for AI software developers as generalist agents[C]//Proceedings of the 13th International Conference on Learning Representations (ICLR 2025). 2025. arXiv: 2407.16741.
  \item WANG X, CHEN Y, YUAN L, et al. Executable code actions elicit better LLM agents[C]//Proceedings of the 41st International Conference on Machine Learning (ICML 2024). 2024. arXiv: 2402.01030.
  \item HONG S, ZHUGE M, CHEN J, et al. MetaGPT: meta programming for a multi-agent collaborative framework[C]//Proceedings of the 12th International Conference on Learning Representations (ICLR 2024). 2024. arXiv: 2308.00352.
  \item QIAN C, LIU W, LIU H, et al. ChatDev: communicative agents for software development[C]//Proceedings of the 62nd Annual Meeting of the Association for Computational Linguistics. 2024: 15174-15186. DOI: 10.18653/v1/2024.acl-long.810.
  \item KAGDI H, COLLARD M L, MALETIC J I. A survey and taxonomy of approaches for mining software repositories in the context of software evolution[J]. Journal of Software Maintenance and Evolution: Research and Practice, 2007, 19(2): 77-131. DOI: 10.1002/smr.344.
  \item ROBBES R, LANZA M. A change-based approach to software evolution[J]. Electronic Notes in Theoretical Computer Science, 2007, 166: 93-109. DOI: 10.1016/j.entcs.2006.06.015.
  \item FALLERI J R, MORANDAT F, BLANC X, et al. Fine-grained and accurate source code differencing[C]//Proceedings of the 29th ACM/IEEE International Conference on Automated Software Engineering (ASE 2014). 2014: 313-324. DOI: 10.1145/2642937.2642982.
  \item GRUND F, CHOWDHURY S A, BRADLEY N C, et al. CodeShovel: constructing method-level source code histories[C]//Proceedings of the 43rd International Conference on Software Engineering (ICSE 2021). 2021: 1510-1522. DOI: 10.1109/ICSE43902.2021.00135.
  \item KO A J, MYERS B A. Designing the Whyline: a debugging interface for asking questions about program behavior[C]//Proceedings of the SIGCHI Conference on Human Factors in Computing Systems (CHI 2004). 2004: 151-158. DOI: 10.1145/985692.985712.
  \item CHENEY J, CHITICARIU L, TAN W C. Provenance in databases: why, how, and where[J]. Foundations and Trends in Databases, 2009, 1(4): 379-474. DOI: 10.1561/1900000006.
  \item PIMENTEL J F, MURTA L, BRAGANHOLO V, et al. noWorkflow: a tool for collecting, analyzing, and managing provenance from Python scripts[J]. Proceedings of the VLDB Endowment, 2017, 10(12): 1841-1844. DOI: 10.14778/3137765.3137789.
  \item PARNIN C, RUGABER S. Resumption strategies for interrupted programming tasks[J]. Software Quality Journal, 2011, 19(1): 5-34. DOI: 10.1007/s11219-010-9104-9.
  \item OPENTELEMETRY. Generative AI semantic conventions[EB/OL]. (2026-10-27). https://\allowbreak opentelemetry.io/\allowbreak docs/\allowbreak specs/\allowbreak semconv/\allowbreak gen-ai/\allowbreak .
  \item LANGFUSE. Observability overview[EB/OL]. (2026-10-27). https://\allowbreak langfuse.com/\allowbreak docs/\allowbreak observability/\allowbreak overview.
  \item 张犬俊, 谢杨, 房春荣, 等. 检索增强生成在软件工程中的应用综述[J]. 软件学报, 2026, 37(3): 1316-1339. DOI: 10.13328/j.cnki.jos.007567.
  \item 汪莹, 字千成, 彭鑫, 等. 基于大语言模型的故障复现测试用例生成方法[J]. 软件学报, 2026, 37(4): 1690-1714. DOI: 10.13328/j.cnki.jos.007474.
\end{tjreflist}}

% ============================================================
\vfill
\tjformsection{六、其他要求}\tjformsecnum{6}

未尽事宜按学院有关规定执行；课题重大调整须经指导教师同意并记录在案。

% ============================================================
% 任务书正文末尾的签字行为四号黑体（审核意见框内的签字行为小四号，见开题/中期报告）
\tjformsignature[\zihao{4}]{指导教师签名}{9em}

% 说明：源稿「八、审核意见」的“指导教师意见／教研室（系）意见／学院意见”三行表，
% 官方任务书表格中不含，故未搬入本文。学院若要求补填，可在上面的签字行之前加入：
%   \begin{tjreviewbox}
%     \tjformlabel{指导教师意见：}
%     \vspace{8\baselineskip}
%     \tjformsignature{指导教师签名}{8em}
%   \end{tjreviewbox}
